% ---------------------------------------------------------------------------
% Capítulo de Introdução.
% Base: capítulos de metodologia, resultados, requisitos e desenvolvimento
% (situação de 17/11/2026), docs/rsl1-protocolo.md, docs/rsl2-protocolo.md e
% docs/referencias-verificadas.md (seção "Introdução (17/11/2026)").
% Os números repetem os dos capítulos; não foram recalculados.
% Escopo de lavagem de dinheiro e corrupção (17/11/2026): fontes em
% docs/lavagem/fontes.md (F01 a F21) e conferência da RSL 1 em
% docs/lavagem/rsl1-lavagem.md.
% Decisão do autor (17/11/2026): a introdução passa a ser capítulo numerado,
% como os demais capítulos do texto, para que as seções não saiam como "0.1".
% ---------------------------------------------------------------------------
\begingroup
\setlength{\emergencystretch}{3em}

\chapter{Introdução}
\label{cap:introducao}

% ===========================================================================
\section{Contextualização}
\label{sec:intro-contexto}

Lavagem de dinheiro é ocultar ou dissimular a natureza, a origem, a
localização, a disposição, a movimentação ou a propriedade de bens, direitos
ou valores provenientes, direta ou indiretamente, de infração penal
\cite[art. 1º]{brasil1998lei9613}. Na redação dada pela Lei nº 12.683/2012,
a lei não lista mais os crimes antecedentes: qualquer infração penal pode
sê-lo \cite[art. 1º]{brasil1998lei9613}. O processo por lavagem independe do
processo da infração antecedente, mas a denúncia precisa de indícios
suficientes de que ela existiu \cite[art. 2º, II e § 1º]{brasil1998lei9613}.

A corrupção é uma dessas possíveis infrações antecedentes (leitura do autor,
já que a lei não lista nenhuma). O Código Penal define a corrupção passiva,
que é solicitar ou receber vantagem indevida em razão da função, e a
corrupção ativa, que é oferecer ou prometer vantagem indevida a funcionário
público \cite[arts. 317 e 333]{brasil1940codigopenal}. Na improbidade
administrativa, é enriquecimento ilícito o agente público adquirir, no
exercício da função e em razão dela, bens decorrentes de vantagem indevida
cujo valor seja desproporcional à evolução do seu patrimônio ou à sua renda,
assegurada a ele a demonstração da licitude da origem
\cite[art. 9º, VII]{brasil1992lei8429}. E a Lei nº 12.846/2013, de natureza
administrativa e civil, inclui entre os atos lesivos à administração pública
usar interposta pessoa física ou jurídica para ocultar os reais interesses ou
a identidade dos beneficiários \cite[art. 5º, III]{brasil2013lei12846}.

Num estudo de casos de lavagem de recursos da corrupção, o Grupo de Ação
Financeira (GAFI) observou
que todos os casos examinados usavam algum tipo de veículo societário,
\textit{trust} ou entidade sem fins lucrativos \cite[§ 47]{fatf2011corruption}.
O GAFI e o Grupo de Egmont listam indicadores de ocultação do beneficiário
final, como um endereço registrado para muitas empresas, administradores que
aparecem em muitas pessoas jurídicas e a troca de administradores ou sócios
logo após a abertura
\cite[anexo E]{fatfegmont2018}. No Brasil, as tipologias de lavagem
publicadas pelo Conselho de Controle de Atividades Financeiras (COAF) incluem
desvio de recursos públicos por empresas contratadas sem licitação e por
convênios com organizações sem fins lucrativos
\cite[tipologias 2.1, 2.7 e 2.8]{coaf2021casos}.

Esses casos mostram que as contratações públicas são uma das portas de
entrada da corrupção e da lavagem, e não a única (leitura do autor). A Lei nº 14.133/2021 estabelece as normas gerais de licitação e
contratação das administrações públicas diretas, autárquicas e fundacionais
da União, dos Estados, do Distrito Federal e dos Municípios
\cite[art. 1º]{brasil2021lei14133}. A mesma lei acrescentou ao Código Penal
um capítulo sobre os crimes em licitações e contratos administrativos, entre
eles a contratação direta ilegal, a frustração do caráter competitivo da
licitação e a fraude em licitação ou contrato
\cite[art. 178]{brasil2021lei14133}. Neste trabalho, as tipologias sobre
contratações sinalizam possíveis infrações antecedentes, e as tipologias de
lavagem sinalizam possível ocultação, pela interposição de pessoas e
empresas, ou enriquecimento incompatível com a renda (decisão do autor).

A Constituição atribui à Polícia Federal as funções de polícia judiciária da União e às
polícias civis, ressalvada a competência da União, as funções de polícia
judiciária e a apuração de infrações penais, exceto as militares
\cite[art. 144, § 1º, IV, e § 4º]{brasil1988constituicao}. O levantamento
de requisitos deste trabalho partiu de uma sugestão do orientador para um
aplicativo OSINT de investigação de corrupção e lavagem de dinheiro na
Polícia Civil do Distrito Federal. Ele tem dois eixos: correlacionar dados
abertos heterogêneos (pessoas, empresas, contratos, bens, redes) e produzir
prova que se sustente em juízo (Seção~\ref{sec:metodologia-requisitos}).

Parte desses dados é aberta. A inteligência de fontes abertas (OSINT, do
inglês \textit{open source intelligence}) é um tipo de inteligência que
aproveita a natureza aberta da coleta, do processamento e da correlação de
dados do ciberespaço para gerar conhecimento \cite{pastorgalindo2020}. Revisões recentes tratam das aplicações de
inteligência artificial (IA) à OSINT \cite{browne2024} e das implicações
éticas, legais e sociais da OSINT com IA \cite{ghioni2024}. Nenhuma das duas
restringe o escopo ao uso por órgãos policiais e de inteligência
(Seção~\ref{sec:metodologia-rsl}).

Ao mesmo tempo, modelos de linguagem de grande porte (LLMs) são aplicados a
tarefas de engenharia de software \cite{hou2024}, inclusive como agentes
\cite{liu2026} e como assistentes dos desenvolvedores \cite{mohamed2026}. Este trabalho
usou agentes de codificação para desenvolver o protótipo
(Seção~\ref{sec:desenv-processo}). Por isso, também reuniu a evidência sobre
esse modo de desenvolver.

% ===========================================================================
\section{Problema}
\label{sec:intro-problema}

Dos dois eixos do levantamento decorrem três condições para uma ferramenta
OSINT de apoio à polícia judiciária na apuração de corrupção e de lavagem de
dinheiro, a partir de contratações públicas, empresas, candidaturas e bens
declarados (decisão do autor). Precisa ler as fontes abertas brasileiras
desses dados. Precisa gerar indícios que o investigador possa conferir. E
precisa preservar cada dado coletado de modo que ele possa servir de prova
no inquérito.

Há dois limites. O primeiro é legal: a infiltração virtual de agentes de
polícia depende de ordem judicial, e é nula a prova obtida sem a observância
dessas regras \cite[art. 10-A]{brasil2013lei12850}. Por isso, a coleta da
ferramenta foi definida como só passiva, para não configurar infiltração
virtual sem ordem judicial (Seção~\ref{sec:requisitos-mercado}). O segundo é
de dados: as informações de contas com o sigilo bancário afastado nos termos
da Lei Complementar nº 105/2001 são prestadas às autoridades, quando
solicitadas, num leiaute definido pelo Banco Central
\cite[art. 1º]{bcb2025in636}, e não são dados abertos. Com dados abertos, a
ferramenta só aponta indícios cadastrais de lavagem, como sócios, endereços,
porte e bens declarados. A movimentação financeira fica numa importação
separada, só para dado recebido por ordem judicial (decisão do autor).

O trabalho investiga se a literatura e as ferramentas existentes atendem a
essas condições e, se não atendem, como um protótipo pode reunir os controles
que faltam. Investiga também como desenvolver esse protótipo com assistentes
e agentes de IA sem perder o controle sobre a qualidade e a segurança do
código. As perguntas de pesquisa de cada revisão estão nas
Tabelas~\ref{tab:qp-rsl1} e~\ref{tab:qp-rsl2}.

% ===========================================================================
\section{Objetivos}
\label{sec:intro-objetivos}

\subsection{Objetivo geral}

Desenvolver um protótipo de buscador OSINT para a polícia judiciária
investigar corrupção e lavagem de dinheiro com dados abertos brasileiros,
tendo as contratações públicas como uma das portas de entrada, com
requisitos fundamentados numa revisão sistemática da literatura, na
legislação e no levantamento com especialista, e com controles de
preservação da prova desde a coleta.

\subsection{Objetivos específicos}

\begin{enumerate}
	\item Identificar e caracterizar, numa revisão sistemática (RSL 1), as
	ferramentas, técnicas e requisitos de sistemas de busca OSINT voltados a
	órgãos policiais e de inteligência, inclusive o uso de IA e os desafios
	legais, éticos e de segurança.
	\item Sintetizar, numa segunda revisão sistemática (RSL 2), as evidências
	empíricas sobre os efeitos de assistentes e agentes baseados em LLMs no
	desenvolvimento de software, nas dimensões de produtividade, qualidade,
	segurança, custo e processo, e derivar delas práticas para o
	desenvolvimento do protótipo.
	\item Confrontar o levantamento de requisitos com a evidência da RSL 1 e
	organizar o resultado em histórias de usuário, agrupadas em épicos e
	priorizadas pelo método MoSCoW.
	\item Comparar o buscador com ferramentas existentes, para identificar o
	que falta a elas e ao protótipo.
	\item Implementar um protótipo com coleta só passiva, cadeia de custódia em
	cada captura, trilha de auditoria, alertas de tipologias de corrupção e
	de lavagem de dinheiro detectáveis com dados abertos, tratados como
	pistas, e IA opcional e local.
	\item Planejar a avaliação do protótipo com usuários policiais.
\end{enumerate}
% PENDENTE: o objetivo 6 depende do capítulo de avaliação, em elaboração. Se
% a avaliação com usuários for executada, trocar "Planejar a avaliação" por
% "Avaliar"; se não for, manter e registrar como trabalho futuro.

% ===========================================================================
\section{Justificativa}
\label{sec:intro-justificativa}

A justificativa vem de duas lacunas, uma na literatura e outra nas
ferramentas.

\textbf{Lacuna na literatura.} A RSL 1 incluiu 57 estudos primários
(Seção~\ref{sec:rsl1-selecao}). Nenhum deles é do Brasil ou de outro país da
América Latina, e nenhum usa dados em português
(Seção~\ref{sec:resultados-rsl1}). Nenhum trata de fraude ou corrupção em
contratações públicas (Seção~\ref{sec:requisitos-mercado}). Nenhum tem como
objeto a lavagem de dinheiro, o beneficiário final ou as empresas de
fachada: as menções à lavagem são de passagem, e o mais próximo da
investigação financeira é o rastreamento de criptoativos
(Seção~\ref{sec:requisitos-novos}). A preservação da
evidência também é exceção: só 7 estudos tratam explicitamente da cadeia de
custódia, e nenhum avalia se a evidência é aceita em juízo
(Seção~\ref{sec:rsl1-qp5}). Só quatro estudos juntam cadeia de custódia e
trilha de auditoria, e nenhum desses quatro tem grafo de vínculos
(Seção~\ref{sec:requisitos-mercado}).

\textbf{Lacuna nas ferramentas.} A comparação com sete ferramentas usou a
documentação oficial de cada uma (Seção~\ref{sec:requisitos-mercado}). As
plataformas de análise de vínculos não mostram, na documentação consultada,
fontes brasileiras integradas nem tipologias de fraude em contratações. As
ferramentas de controle da Controladoria-Geral da União (CGU) cobrem o
domínio, mas são ferramentas de auditoria, e a documentação consultada não
descreve preservação da evidência com \textit{hash} nem trilha de auditoria
das consultas. Na ferramenta de captura de evidências, a documentação
consultada não mostra grafo de vínculos nem tipologias. Segundo a
documentação consultada, nenhuma ferramenta comparada reúne esses recursos
numa só ferramenta para a polícia judiciária. A comparação usou critérios
de contratações públicas; as tipologias de lavagem não fizeram parte dela.

\textbf{Desenvolvimento com IA.} A RSL 2 mostra que o ganho de produtividade
com IA é condicional e que o código gerado não é seguro por padrão
(Seções~\ref{sec:resultados-rsl2-qp1} e~\ref{sec:resultados-rsl2-qp3-qp4}).
Um protótipo para uso policial, desenvolvido com agentes de IA, precisa de
práticas que controlem esses riscos. A RSL 2 também não encontrou estudo
sobre o desenvolvimento individual de um produto completo em contexto
acadêmico (Seção~\ref{sec:resultados-rsl2-qp5}).

% ===========================================================================
\section{Contribuições}
\label{sec:intro-contribuicoes}

As contribuições do trabalho são:
\begin{itemize}
	\item uma revisão sistemática sobre ferramentas OSINT para órgãos policiais
	e de inteligência, com a rastreabilidade de cada achado para os requisitos
	do buscador (Seção~\ref{sec:resultados-rsl1});
	\item uma revisão sistemática sobre os efeitos de assistentes e agentes de
	IA no desenvolvimento de software, com uma lista de práticas candidatas
	para o desenvolvimento (Seção~\ref{sec:resultados-rsl2});
	\item um \textit{backlog} priorizado, em que cada história aponta o
	requisito que atende e a evidência que a sustenta, inclusive as
	tipologias de lavagem de dinheiro detectáveis com dados abertos e os
	limites desses dados, e uma comparação com ferramentas existentes
	(Capítulo~\ref{cap:requisitos});
	\item um protótipo de código aberto, sob a licença MIT, que reúne coleta
	passiva, cadeia de custódia, trilha de auditoria, alertas explicados e IA
	opcional e local, com o modelo de ameaças ligado aos testes automatizados
	(Capítulo~\ref{cap:desenvolvimento});
	\item o registro do processo de desenvolvimento com agentes de IA, com
	modelo, versão, duração e \textit{tokens} de cada execução
	(Seção~\ref{sec:desenv-processo}).
\end{itemize}

Os resultados das duas revisões são provisórios. A leitura do texto completo,
o \textit{snowballing}, a avaliação de qualidade, a extração e a síntese foram
feitos pela IA como terceira opinião, com regras aceitas pelo autor (R1). O
orientador (R2) ainda confirma essas regras, e a avaliação independente dos
dois revisores na etapa 2 ainda não foi feita
(Seções~\ref{sec:resultados-rsl1} e~\ref{sec:resultados-rsl2}). Os testes
automatizados do protótipo usam só dados sintéticos, e as fontes reais
passaram só por um teste de fumaça. O protótipo ainda não foi avaliado com
policiais
(Capítulo~\ref{cap:desenvolvimento} e Seção~\ref{sec:desenv-limitacoes}).
% PENDENTE: atualizar este parágrafo quando R2 confirmar as regras e a
% leitura independente de R1 e R2 for feita.

% ===========================================================================
\section{Organização do trabalho}
\label{sec:intro-organizacao}

O Capítulo~\ref{cap:metodologia} descreve a classificação da pesquisa, as
etapas do trabalho, o método das duas revisões sistemáticas, a forma de
registro que as torna verificáveis e o levantamento de requisitos. Inclui a
declaração do uso de IA nas revisões.

O Capítulo~\ref{cap:resultados-rsl} apresenta os resultados das duas
revisões: a seleção dos estudos, as características, a qualidade e a síntese
de cada questão de pesquisa. A RSL 1 termina na rastreabilidade dos achados
para os requisitos do buscador, e a RSL 2, nas práticas candidatas para o
desenvolvimento.

O Capítulo~\ref{cap:requisitos} junta os requisitos do levantamento e da RSL
1, compara o buscador com ferramentas existentes, escreve os requisitos como
histórias de usuário e os prioriza.

O Capítulo~\ref{cap:desenvolvimento} descreve o protótipo: a arquitetura, os
casos de uso, o modelo de dados, os fluxos, a interface de programação, as
decisões de arquitetura, o modelo de ameaças, a verificação, o processo de
desenvolvimento com IA e as limitações.

% PENDENTE: quando a avaliação for executada, incluir os resultados aqui.
O Capítulo~\ref{cap:avaliacao} apresenta o planejamento da avaliação do protótipo
com usuários policiais.

O Capítulo~\ref{cap:consideracoes} retoma os objetivos, os principais
achados, as contribuições, as limitações e os trabalhos futuros.

\endgroup
